\documentclass[10pt,a4paper]{amsart}
\usepackage[margin=19mm]{geometry}
\usepackage{amsmath,amssymb,mathtools}
\usepackage[scr=rsfs,cal=euler]{mathalfa}
\usepackage{xfrac}
\usepackage[unicode,hidelinks,bookmarks=false]{hyperref}
\newcommand{\ie}{i.e.\ }
\newcommand{\cf}{cf.\ }
\newcommand{\resp}{respectively\ }
\newcommand{\Subsection}{\subsection}
\providecommand{\nobreakdash}{\nobreak\hbox{-}}
\setlength{\emergencystretch}{2em}
\multlinegap0pt
\usepackage{mathtools}
\usepackage{enumitem,booktabs,tikz,fullpage}
\allowdisplaybreaks[4]
\newtheorem{theorem}{Theorem}%[section]
\newtheorem{lemma}{Lemma}[section]
\newtheorem{proposition}{Proposition}[section]
\newtheorem{claim}{Claim}[section]
\newtheorem*{remark}{Remark}
\newtheorem{conjecture}{Conjecture}
\DeclareMathOperator{\mc}{mc}
\begin{document}
\title[Proving a conjecture concerning chromatic number, size and l]{Proving a conjecture concerning chromatic number, size and least eigenvalue}
\author{Guangzhou 510631, P.R. China; Leyou Xu, Bo Zhou}
\date{}
\maketitle
\noindent\emph{Source and licence.} Proving a conjecture concerning chromatic number, size and least eigenvalue. Version arXiv:2608.03271v1 (CC BY 4.0). This material is distributed under the Creative Commons Attribution 4.0 International licence, http://creativecommons.org/licenses/by/4.0/, as recorded in the arXiv OAI licence metadata for this exact version and shown on the abstract page https://arxiv.org/abs/2608.03271v1. Full rights notice: licenses/arxiv-2608.03271-CC-BY-4.0.txt.\par\smallskip
\noindent\textit{Editorial scope.} Counted unit: the original \texttt{theorem} environment \texttt{X} at source lines 86--89 together with its complete proof at lines 232--531; the delivered document keeps source lines 51--532 so that every referenced label and every citation resolves inside it. Native editable source is the paper's own concatenated TeX; the AMS wrapper is editorial formatting only. No publisher class, upstream script or unrelated artwork is redistributed.
\section{Introduction}

Throughout this article, all graphs  are finite and simple. Let $G$ be a graph with vertex set $V(G)$ and edge set $E(G)$. The size of $G$ is $e(G)=|E(G)|$. 
A nonempty graph is a graph with size at least one. For two graphs $G$ and $H$, $G\cup H$ denotes the disjoint union of $G$ and $H$, and $G\vee H$ denotes the join of $G$ and $H$. The notations $K_n$ and $K_{a,b}$ denote the complete $n$-vertex graph and the complete bipartite graph with partite sizes $a$ and $b$. By $nK_1$ we denote the graph consisting of $n$ isolated vertices. By convention, $G\cup 0K_1=G$ for any graph $G$.
The chromatic number of a graph $G$, denoted by $\chi(G)$, is the minimum number of colors needed to color the vertices of $G$ such that no two adjacent vertices share the same color. 

Let $G$ be a graph of order $n$. 
Let $A(G)$ be the adjacency matrix of $G$. The eigenvalues of $A(G)$ are called the eigenvalues of  $G$, which are ordered as $\lambda_1(G)\ge \cdots\ge \lambda_n(G)$.  The least eigenvalue $\lambda_n(G)$ of $G$ is  denoted $\lambda (G)$. The largest eigenvalue $\lambda_1(G)$ is also known as the spectral radius of $G$. By Perron-Frobenius theorem, $-\lambda (G)\le \lambda_1(G)$.

For a graph parameter $f(G)$, when no confusion arises, we omit $G$ and simply write $f$. 

For a graph $G$ of order $n$, Tang and Elphick proved that if $3\le \chi\le n-1$, 
\begin{equation}\label{old}
\chi\le \frac{n}{2}+1+\lambda
 + \sqrt{\left(\frac{n}{2}+1+\lambda\right)^2-4(\lambda+1)\left(\lambda+\frac{n}{2}\right)}
\end{equation}
with equality if and only if $G\cong (K_{\frac{\chi}{2}}\cup \frac{n-\chi}{2}K_1)\vee (K_{\frac{\chi}{2}}\cup \frac{n-\chi}{2}K_1)$, where both $n$ and $\chi$ are even. For  $3\le \chi\le \frac{n}{2}$, this was already proved by Fan et al.~\cite{FYW} and they conjectured the range $3\le \chi\le \frac{n}{2}$ can be  extended to  $3\le \chi\le n-1$.  
Tang and Elphick observed that $\frac{n}{2}$ in \eqref{old} is closely related to the bound $\lambda\ge -\frac{n}{2}$ \cite{Con,Pow}. Let $m$ be the size of $G$. Wu and Elphick \cite{WuElphick} proved that
$\chi(\chi-1)\le (\lambda_1+1)\lambda_1\le 2m$. By analogy with \eqref{old} to replace $\frac{n}{2}$ with $m$, $\chi$ with $\chi(\chi-1)$ and $\lambda$ with $-\lambda^2$, Tang and Elphick \cite{TangElphick} proposed the following conjecture. 

\begin{conjecture} \label{con} \cite[Conjecture~8]{TangElphick}
For any nonempty graph $G$ with size $m$,
\begin{equation}\label{eq:main}
\chi(\chi-1) \le m+1-\lambda^2
 + \sqrt{(m+1-\lambda^2)^2-4(\lambda^2-1)(\lambda^2-m)}.
\end{equation}
\end{conjecture}

%Tang and Elphick reported their verification of this conjecture  for several graph databases and bipartite graphs.  
Tang and Elphick observed that the conjecture is immediate for bipartite graphs. They  verified it for all graphs of order at most nine and for the graphs of order at most one hundred in the Wolfram Mathematica database.
They also compared \eqref{eq:main} and \eqref{old} and found that the bound in \eqref{eq:main} typically performs better than the bound in \eqref{old}.


In this paper, we show that Conjecture \ref{con} is true and characterize the equality completely.


\begin{theorem}\label{X}
For any nonempty graph $G$ with size $m$, \eqref{eq:main} follows, and  
equality holds  if and only if $G\cong K_\chi\cup (n-\chi)K_1$ or when $\chi=2$ and $G\cong K_{a,b}\cup (n-a-b)K_1$ for some  $a,b\ge 1$, where $n$ is the order of $G$.
\end{theorem}



\section{Preliminaries}

Let $G$ be a graph. For $\emptyset\ne S\subseteq V(G)$, let $G[S]$ denote the subgraph of $G$ induced by $S$. If $S\subset V(G)$, then $G-S=G[V(G)\setminus S]$.
For $u\in V(G)$, we write $G-u$ for $G-\{u\}$. 
For $E_1\subseteq E(G)$, $G-E_1$ denotes the spanning subgraph of $G$ obtained by deleting edges of $E_1$. In particular, if $E_1=\{e\}$, we simply write $G-e$ for $G-\{e\}$.


For $v\in V(G)$, let $d_G(v)$ denote the degree of $v$ in $G$. If $S\subseteq V(G)$, we  simply write $d_S(v)$ for the number of neighbors of $v$ in $S$. %$d_{G}(v)\cap S$. 
For disjoint subsets $X,Y\subset V(G)$, we denote by $e(X,Y)$ the number of edges between vertices of $X$ and vertices of $Y$.

 
A clique in a graph is a set of vertices that are pairwise adjacent. The maximum size of a clique of $G$ is called the clique number of $G$, denoted by $\omega(G)$. Evidently, $\omega(G)\le \chi(G)$. 

Let $r$ be a positive integer. An $r$-chromatic graph is a graph whose chromatic number is exactly $r$. Every $r$-chromatic graph has size at least ${r\choose 2}$. 
A graph $G$ is $r$-vertex-critical if it is $r$-chromatic and  
$\chi(G-v)\le r-1$ for every $v\in V(G)$ \cite{Jen}.  
Every $r$-chromatic graph contains an induced
$r$-vertex-critical subgraph. 
Note that  for $r\ge 4$, the set of  $r$-vertex-critical graphs and the  set of $r$-critical graphs (those $G$ with $\chi (G)=r$ and  $\chi(G-t)\le r-1$ for every $t\in V(G)\cup E(G)$) are different.




%: repeatedly delete a vertex whenever its deletion leaves chromatic
%number $r$. 
%
%Moreover, every $r$-critical graph $F$ has minimum degree at least $r-1$.
%
%otherwise, an $(r-1)$-coloring of $F-v$ could be extended to a vertex $v$ of degree at most $r-2$. 
%This is the following fact, see \cite{CGSZ,KW}. 

\begin{lemma}\label{edge} \cite{Jen}
An $r$-vertex-critical graph has minimum degree at least $r-1$.  
\end{lemma}


%As a further consequence, we obtain the following.

\begin{lemma}\label{order}
For a noncomplete  $r$-vertex-critical graph $G$ with  $r\ge 3$, $|V(G)|\ge r+2$. 
\end{lemma}
\begin{proof}  Suppose that  $|V(G)|< r+2$. Then, as $G$ is not complete, $|V(G)|=r+1$, and there are two vertices $u$ and $v$ that  are not adjacent. By Lemma~\ref{edge}, both $u$ and $v$ are adjacent to all other vertices of $G$. Moreover, if there exists another vertex pair $\{w,z\}$ that are not adjacent, then we may color $\{w,z\}$ with one color, $\{u,v\}$ with another color and the remaining $r-3$ vertices with distinct colors. This gives an $(r-1)$-coloring, a contradiction. So $G\cong K_{r+1}-uv$, which is not $r$-vertex-critical, a contradiction. 
\end{proof}



We shall need the following structural lemma, which guarantees the existence of a clique that is critical in the sense of vertex deletion. 

\begin{lemma}\label{lem:small-extremal-coloring}
Let $G$ be an $r$-chromatic graph with $r\ge 2$ and $e(G) \le \binom{r}{2}+r-2$.
Then $G$ contains an $r$-clique $C$ such that, for every $u\in C$,
$\chi(G-u)\le r-1$.
\end{lemma}

\begin{proof}
If $r=2$, then $e(G)=1$, and $G\cong K_2\cup tK_1$ for some $t$, so the result follows. 

Assume that $r\ge 3$. 
Let $F$ be an induced $r$-vertex-critical subgraph of $G$. 
If $F\ncong K_r$, then by Lemma~\ref{order}, we have $|V(F)|\ge r+2$ and so by Lemma~\ref{edge}, $e(G)\ge e(F)\ge \frac{(r-1)(r+2)}{2}=\binom{r}{2}+r-1$, a contradiction. Hence $F\cong K_r$. Then $C=V(F)$ is an $r$-clique of $G$.

Let $u$ be an arbitrary vertex of $C$.
First, color all vertices of $C\setminus\{u\}$ with  colors $1,\dots, r-1$. Let $F_0=G-C$. As $e(G) \le \binom{r}{2}+r-2$, for each $v\in V(F_0)$, let 
\[
L(v)=\{i\in \{1,\dots, r-1\}: \text{$i$ is not used on neighbors of $v$ in $C\setminus\{u\}$}\}. 
\]
As $e(G)\le \binom{r}{2}+r-2$, there are at most $r-2$ edges of $G$ with one end outside $C$, so for each $v\in V(F_0)$,
$d_{C\setminus\{u\}}(v)+d_{F_0}(v)\le r-2$, implying that 
$|L(v)|\ge (r-1)-d_{C\setminus\{u\}}(v)\ge d_{F_0}(v)+1$.
We then color the vertices of $F_0$ greedily in any order. When a vertex $v$ is colored, at most $d_{F_0}(v)$ colors are forbidden by its already colored neighbors, so an available color remains
in $L(v)$. This gives an $(r-1)$-coloring of $G-u$, so $\chi(G-u)\le r-1$.
\end{proof}

Dirac  established a lower bound  on the size of a noncomplete $r$-vertex-critical graph for $r\ge 4$, see also \cite{BernshteynKostochka}.

\begin{theorem}\label{thm:dirac} \cite[Theorem 15]{Dirac}
Let $r\ge 4$, and let $G$ be a noncomplete $r$-vertex-critical graph. Then
$2e(G)\ge (r-1)|V(G)|+r-3$.
\end{theorem}

%Given a connected graph $G$, we have by Perron-Frobenius theorem that there is a unique positive unit eigenvector corresponding to $\lambda_1$.%, which we call the Perron vector of $G$. 
Given a graph $G$ that is not necessarily connected, we have by Perron-Frobenius theorem applied to each component, there is a nonnegative unit eigenvector corresponding to $\lambda_1$.
This implies the following well known lemma.

\begin{lemma}\label{subgraph}
If $F$ is a subgraph of a graph $G$, then $\lambda_1(F)\le \lambda_1(G)$.
\end{lemma}

The bound in the following lemma is known for (connected) graphs \cite{Powers}  and the equality case for connected bipartite graphs is also known \cite{HS}. However, for completeness, we include a proof here.


\begin{lemma} \label{BIS}
For a graph $G$ of size $m\ge 1$, $\lambda^2\le m$ with equality if and only if $G$ is a complete bipartite graph with possibly isolated vertices.
\end{lemma}

\begin{proof} 
Let $n=|V(G)|$. From $\sum_{i=1}^n\lambda_i^2=2m$, one has $\lambda^2+\lambda_1^2\le 2m$. As $-\lambda \le \lambda_1$, it follows that $\lambda^2\le m$ \cite{Powers}.

If $G$ is a complete bipartite graph, then it is trivial that  $\lambda^2=\lambda_1^2=m$.

Suppose that $\lambda^2=m$. Then all inequalities above are equalities, so $-\lambda=\lambda_1$ and $\lambda_i=0$ for $1<i<n$. As $G$ has exactly one positive eigenvalue, $G$ has only one non-trivial connected component, say $F$. By Perron-Frobenius theorem, $F$ is bipartite (otherwise, the index of imprimitivity of $A(F)$ is $1$, so $\lambda$ cannot be an eigenvalue of $F$ and so $G$). 
%As $F$ is bipartite and has exactly one positive eigenvalue, $F$ is a complete bipartite graph \cite{Pe}. 
Then $F$ is a complete bipartite graph as otherwise $F$ contains a $4$-vertex induced path, so Cauchy's interlacing theorem (see \cite[Theorem 0.10]{CDS}) implies that  $F$ has two neigative eigenvalues, a contradiction.
\end{proof}

Given a nonempty graph $G$, let $\mathbf{x}$ be an eigenvector associated with $\lambda$. %Then $\mathbf{x}$ is orthogonal to the Perron vector of $G$. %not true, as $G$ may not be connected
Then $\mathbf{x}$ has both positive and negative entries. Let $X = \{v : x_v \ge 0\}$ and $Y = \{v : x_v < 0\}$. Let $H$ be the spanning subgraph of $G$ with $E(H)$ to be the set of edges of $G$ between vertices $X$ and vertices of $Y$. Then $H$ is a bipartite graph with bipartition $(X, Y)$, which we call the bipartite graph of $G$ determined by $\mathbf{x}$.


%Let $G$ be a graph of order at least two. For $S\subset V(G)$, $\partial_G(S)$ denotes the set of edges of $G$ with one end vertex in $S$ and the other vertex outside $S$. 
%%A cut of $G$ is a partition of $V(G)$ into two nonempty parts $A$ and $B$. An edge is said to be cut by $(A,B)$ if it has one endpoint in $A$ and the other in $B$, otherwise it is uncut. We also call the set of edges crossing between $A$ and $B$ a cut.  Let $e(A,B)$ denote the number of edges crossing between $A$ and $B$. 
%We define $\operatorname{mc}(G) = \max_{S\subset V(G)} |\partial_G(S)|$.

\begin{lemma}\label{maxcut}
For any nonempty graph $G$ and its  bipartite graph $H$ determined by some eigenvector  $\mathbf{x}$ associated with $\lambda$, $\lambda^2 \le e(H)$.
\end{lemma}

\begin{proof}
%Let $X = \{v : x_v \ge 0\}$ and $Y= \{v : x_v < 0\}$. 
Let $\mathbf{z}$ be a vector with  $z_v = |x_v|$ for $v\in V(G)$. From Rayleigh's principle,
\[
\lambda \mathbf{x}^\top\mathbf{x}=\mathbf{x}^\top A(G)\mathbf{x}=2\sum_{uv\in E(G)}x_ux_v\ge 2\sum_{uv\in E(H)}x_ux_v=\mathbf{x}^\top A(H)\mathbf{x}\ge \lambda(H)\mathbf{x}^\top\mathbf{x},
\]
so $\lambda\ge \lambda(H)$, implies that $\lambda^2\le \lambda^2(H)$. Now the result follows from Lemma \ref{BIS}.
\end{proof}
%\[
%-\lambda z_v=-\sum_{vw\in E(G)\setminus E(H)}z_w+\sum_{vw\in E(H)} z_w \le \sum_{vw\in E(H)} z_w
%\]
%and so 
%\[
%-\lambda\mathbf{z}^\top\mathbf{z}=-\lambda\sum_{v\in V(G)}z_v^2\le \sum_{v\in V(G)}z_v\sum_{vw\in E(H)} z_w=\mathbf{z}^\top A(H)\mathbf{z}\le \lambda_1(H)\mathbf{z}^\top\mathbf{z}.
%\]
%As $\mathbf{z}\neq \mathbf{0}$, it follows that $-\lambda\le \lambda_1(H)$. As $H$ is bipartite, $\lambda_1(H)\le \sqrt{e(H)}$. It then follows that $\lambda^2\le e(H)$.



\section{Proof of Theorem~\ref{X}}

We are now ready to prove Theorem~\ref{X}. To prove the inequality, we first transform \eqref{eq:main} into an equivalent form, and then use a case distinction based on the relationship between the clique number and the chromatic number  and the properties of the $r$-vertex-critical graphs.


\begin{proof}[Proof of Theorem~\ref{X}]
%To begin with, we verify that equality is attained by the two given graph families. If
%$G\cong K_\chi\cup (n-\chi)K_1$, then $m=\binom \chi2$ and $\lambda=-1$ and so \eqref{eq:main} becomes an equality.
%If $G\cong K_{a,b}\cup (n-a-b)K_1$ with $a,b\ge 1$, then $m=ab$ and $\lambda=-\sqrt{ab}$, so \eqref{eq:main} becomes an equality.
%
%We now prove the inequality for an arbitrary graph and show that there are no other equality cases. 

 
As $G$ is nonempty, we have $\chi\ge 2$. Also, as $G$ is nonempty,  we have by Cauchy's interlacing theorem that $\lambda=\lambda (G)\le \lambda(K_2)= -1$, so $\lambda^2\ge 1$. By Lemma \ref{BIS}, $\lambda^2\le m$. Thus
 $(\lambda^2-1)(\lambda^2-m)\le 0$ and $m+1-\lambda^2\ge 1$. 

%Denote by $n$ the order of $G$. %Theorem 1.1 states $n$ is the order of $G$

Suppose that $\chi=2$. Then \[
m+1-\lambda^2+\sqrt{(m+1-\lambda^2)^2-4(\lambda^2-1)(\lambda^2-m)}\ge 2(m+1-\lambda^2)\ge 2
\]
with equalities if and only if $\lambda^2=m$. So \eqref{eq:main} holds, and  by Lemma \ref{BIS}, 
\eqref{eq:main} is equality if and only if $G\cong K_{a,b}\cup (n-a-b)K_1$ with $a,b\ge 1$.

Suppose that $\chi\ge 3$. 
Note that $G$ has at least $\binom{\chi}{2}$ edges. Let $s = m-\binom{\chi}{2}$. Then $s\ge 0$. Let 
\[
f(x)=x^2 - 2(m+1-\lambda^2)x + 4(\lambda^2-1)(\lambda^2-m).
\]
%By Cauchy's interlacing that $\lambda=\lambda (G)\le \min\{\lambda(K_{1,2}), \lambda(K_3)\}< -1$, so $\lambda^2> 1$. By Lemma \ref{BIS}, $\lambda^2< m$. Thus $(\lambda^2-1)(\lambda^2-m)< 0$.  %no need for strict inequality
Recall that  $(\lambda^2-1)(\lambda^2-m)\le 0$ and $m+1-\lambda^2\ge 1$. 
Then $f(x)=0$ has two roots $r_1$ and $r_2$ with $r_1< r_2$ and $r_1\le 0$. So when $x>0$, $x\in (r_1,r_2] \iff f(x)\le 0$ with equality if and only if $x=r_2$. 
Observe  that  \eqref{eq:main} becomes $\chi(\chi-1)\le r_2$. 
As $\chi(\chi-1)>0$, \eqref{eq:main} is equivalent to $f(\chi(\chi-1))\le 0$, i.e.,
\[
(\chi(\chi-1))^2-2(m+1-\lambda^2)\chi(\chi-1)+4(\lambda^2-1)(\lambda^2-m)\le 0.
\]
As $m=s+{\chi \choose 2}$, the above inequality is equivalent to 
\[
\lambda^4-(s+1)\lambda^2-s\left({\chi \choose 2}-1\right)\le 0,
\] 
that is, 
\[
\lambda^2 \le U_\chi(s) := \frac{s+1+\sqrt{(s-1)^2 + 2\chi(\chi-1)s}}{2}.
\]
Moreover, equality holds in \eqref{eq:main}  if and only if $\lambda^2= U_\chi(s)$. Therefore it suffices to prove that $\lambda^2 \le U_\chi(s)$ with equality if and only if $G\cong K_\chi \cup (n-\chi)K_1$.


 

Let $H$ be the bipartite graph of $G$ determined by some eigenvector associated with $\lambda$. Let $(X,Y)$ be its bipartition.

\noindent {\bf Case 1.}
 $\omega(G)=\chi$.



Let $C$ be a fixed $\chi$-clique of $G$.

If $s=0$, then $G\cong K_\chi\cup (n-\chi) K_1$, as desired. Suppose that $s\ge 1$.   
It suffices to show that $\lambda^2 <U_\chi(s)$.

We first show an upper bound on $-\lambda$. Let $\tau = \frac{1+\sqrt{1+4s}}{2}$. Then $\tau(\tau-1)=s$.
\begin{claim}\label{c1}
$-\lambda< \tau$. %with equality if and only if $s=0$.
%$\lambda(\lambda+1) \le s$ with equality if and only if $s=0$. 
\end{claim}
\begin{proof}
Let $E_1=E(G[C])$. Let $M=A(G-E_1)$ and $A_0=A(G)-M$. 
Note that $A_0 + \tau I = J_\chi + D$, where $J_\chi$ is a matrix with a $\chi\times \chi$ principal submatrix of all  ones (corresponding to the vertices of $C$) and zeros elsewhere, and $D$ is a diagonal matrix with \[
D_{vv}=\begin{cases}
\tau-1, & v\in C,\\
\tau, & v\notin C.
\end{cases}
\]

We show that $D+M$ is positive semidefinite.
Let $\mathbf{x}\in  \mathbb{R}^{n}$ and $\mathbf{z}$ be a vector with $z_v=|x_v|$ for any $v\in V(G)$. Let $V^+ = \{v : x_v > 0\}$, $V^- = \{v : x_v < 0\}$ and $F$ be the spanning bipartite subgraph of $G-E(C)$ induced by the edges between $V^+$ and $V^-$.
Then \begin{align*}
\mathbf{x}^\top (D+M)\mathbf{x}&=\sum_{v\in V(G)}D_{vv}x_v^2+2\sum_{uv\in E(G-E(C))}x_ux_v\\
&\ge \sum_{v\in V(G)}D_{vv}z_v^2-2\sum_{uv\in E(F)}z_uz_v=\mathbf{z}^\top (D-A(F))\mathbf{z}.
\end{align*}
Let $W = D^{-1/2}A(F)D^{-1/2}$.
Let $\mathbf{y}$ be a nonnegative unit eigenvector corresponding to $\rho(W)$, where $\rho(W)$ is the spectral radius of $W$. Then \[
\rho(W)=2\sum_{uv\in E(F)}w_{uv}y_uy_v.
\]
By Cauchy-Schwarz inequality, we have \begin{align*}
\rho(W)^2 &\le \left(\sum_{uv\in E(F)}w_{uv}^2\right)\left( \sum_{uv\in E(F)}4y_u^2y_v^2\right)\\
&\le 4\left(\sum_{uv\in E(F)}w_{uv}^2\right)\left(\sum_{u\in V^+}y_u^2\right)\left(\sum_{u\in V^-}y_u^2\right)\le \sum_{uv\in E(F)}w_{uv}^2,
\end{align*}
where the second inequality follows as $F$ is bipartite, and the third inequality follows as $\sum_{u\in V^+}y_u^2+\sum_{u\in V^-}y_u^2\le 1$.
As every edge in $F$ either joins $C$ to $V(G)\setminus C$ or lies in $V(G)\setminus C$, we have $w_{uv}^2=\frac{1}{\tau(\tau-1)}=\frac{1}{s}$ in the former case and $w_{uv}^2=\frac{1}{\tau^2}<\frac{1}{s}$. Then $\sum_{uv\in E(F)}w_{uv}^2\le 1$. This shows that $\rho(W)\le 1$, implying that $I-W$ is positive semidefinite.
Note that $D-A(F)=D^{1/2}(I-W)D^{1/2}$. Thus $D-A(F)$ is positive semidefinite, and hence $D+M$ is positive semidefinite, as desired. 

As both  $J_{\chi}$ and $D+M$ are positive semidefinite, 
 $A(G)+\tau I=J_{\chi}+D+M$ is positive semidefinite, so  $-\lambda\le \tau$. 

Suppose that $-\lambda=\tau$. From the preceding argument, there exists a non-zero vector $\mathbf{x}$ such that $\sum_{v\in C}x_v=0$, $\mathbf{z}^\top (D-A(F))\mathbf{z}=0$ and $\rho(W)=1$, where $\mathbf{z}$ is the vector corresponding to $\mathbf{x}$ defined above. From $\rho(W)=1$, we have $F$ contains all $s$ edges outside $C$, each of these edges connects a vertex of $C$ and a vertex outside $C$. Consequently, $W=\tfrac{1}{\sqrt{s}}A(F)$ and $\lambda_1(F)^2=s=e(F)$, implying that $F$ has a non-trivial complete bipartite component with $s$ edges. So the vertices of $C$ in $F$ are all contained in either $V^+$ or $V^-$, contradicting $\sum_{v\in C}x_v=0$. So $-\lambda<\tau$.
\end{proof}

By Claim~\ref{c1},
\[
\lambda^2<\tau^2=\left(\frac{1+\sqrt{1+4s}}2\right)^2=s+\frac12+\frac12\sqrt{1+4s}.
\]
To prove $\lambda^2<U_\chi(s)$, it suffices to show that $s+\frac12+\frac12\sqrt{1+4s}\le U_\chi(s)$, or equivalently, %$s+\sqrt{1+4s}\le\sqrt{(s-1)^2+2k(k-1)s}$, which is in turn equivalent to 
$\sqrt{1+4s}\le \chi(\chi-1)-3$, and this is further equivalent to \[
s\le\frac{(\chi(\chi-1)-3)^2-1}{4}=:A_\chi.
\]
So the result holds if $s\le A_\chi$. 

Suppose next that $s>A_\chi$.
 Let $a=|X\cap C|$ and $b=|Y\cap C|$. Then $a+b=\chi$ and 
\[
e(G[X])+e(G[Y])\ge {a\choose 2}+{b\choose 2}=a^2-\chi a+\frac{1}{2}(\chi^2-\chi)
\ge \left\lfloor\frac{(\chi-1)^2}{4}\right\rfloor.
\]
So 
Lemma~\ref{maxcut} gives
$$\lambda^2\le e(H)=m-(e(G[X])+e(G[Y]))\le m-\left\lfloor\frac{(\chi-1)^2}{4}\right\rfloor=s+\binom{\chi}{2}-\left\lfloor\frac{(\chi-1)^2}{4}\right\rfloor.$$
To prove $\lambda^2 <U_\chi(s)$, it suffices to show that
 $s+\binom{\chi}{2}-\left\lfloor\frac{(\chi-1)^2}{4}\right\rfloor< U_\chi(s)$, or equivalently, 
\[
s> \frac{\left(\chi(\chi-1)-2\left\lfloor\frac{(\chi-1)^2}{4}\right\rfloor\right)\left(\chi(\chi-1)-2\left\lfloor\frac{(\chi-1)^2}{4}\right\rfloor-2\right)}{4\left\lfloor\frac{(\chi-1)^2}{4}\right\rfloor}=:R_\chi.
\]
 If $\chi=2h\ge4$, then
\[
A_\chi-R_\chi=2(2h^2-1)(h^2-h-1)\ge0,
\]
and if $\chi=2h+1\ge3$, then
\[
A_\chi-R_\chi=
\frac{(h-1)(h+1)(4h^3+4h^2-2h-1)}h\ge0.
\]
Thus  $s>A_\chi\ge  R_\chi$, as desired. 


\noindent {\bf Case 2.}  $\omega(G)<\chi$. 

It suffices to show that $\lambda^2 <U_\chi(s)$.

\noindent {\bf Case 2.1.}  $\chi=3$. 

Note that  $G$ is triangle free and non-bipartite. As $H\ne G$, Lemma~\ref{maxcut} gives $\lambda^2\le m-1=s+2$. It suffices to show  $s+2<U_3(s)$, i.e., 
\[
s+2<\frac{1}{2}\left(s+1+\sqrt{(s-1)^2+12s} \right),
\] 
 or equivalently, $s> 2$. In fact, a triangle-free non-bipartite graph contains an odd cycle of length at least five, so $m\ge5$ and
$s=m-3$. If $m\ge 6$, then $s>2$, as desired. If $m=5$, then 
$G\cong C_5\cup (n-5)K_1$, and by a direct calculation, $\lambda=-\frac{1+\sqrt5}{2}$, and hence
$\lambda^2=\frac{3+\sqrt5}{2}<4=U_3(2)$. 

\noindent {\bf Case 2.2.}  $\chi\ge4$. 

Let
$\beta=m-e(H)$. %=e(G[X])+e(G[Y])
%Then $\beta(G)$ is the minimum number of edges left inside the two parts of a cut. 
Let $B_\chi=\left\lfloor\frac{(\chi-1)^2}{4}\right\rfloor
 +\left\lfloor\frac{\chi}{2}\right\rfloor-1$.

\begin{claim}\label{lem:maxcut-deficiency}
 $\beta\ge B_\chi$.
\end{claim}

\begin{proof}
Let $a=\chi(G[X])$ and $b=\chi(G[Y])$. 
As $G[X]$ and $G[Y]$ can be colored with disjoint sets of colors, $a+b\ge \chi$. Then
\[
\beta\ge\binom a2+\binom b2.
\]

Suppose to the contrary that $\beta<B_\chi$. As $\beta$ is an integer,
\[
\beta\le \left\lfloor\frac{(\chi-1)^2}{4}\right\rfloor
 +\left\lfloor\frac{\chi}{2}\right\rfloor-2\le \begin{cases}
h^2-2 &\text{if }\chi=2h,\\
h^2+h-2 &\text{if }\chi=2h+1.
\end{cases}
\]
If $a+b\ge \chi+1$, then
\[
\binom a2+\binom b2\ge
\begin{cases}
h^2 &\text{if }\chi=2h,\\
h^2+h &\text{if }\chi=2h+1,
\end{cases}
\]
a contradiction. So $a+b=\chi$.

Assume that $a\le b$. 
If $\chi=2h$, write $a=h-r$ and $b=h+r$, where $r\ge0$. Then
$\binom a2+\binom b2=\left\lfloor\frac{(\chi-1)^2}{4}\right\rfloor+r^2$, 
and hence
\[
\beta-\binom a2-\binom b2
 \le h-2-r^2\le h-r-2=a-2.
\]
If $\chi=2h+1$, write $a=h-r$ and $b=h+1+r$. Then
$\binom a2+\binom b2=\left\lfloor\frac{(\chi-1)^2}{4}\right\rfloor+r(r+1)$,
and hence
\[
\beta-\binom a2-\binom b2
 \le h-2-r(r+1)\le h-r-2=a-2.
\]
So in both cases, $\beta-\binom a2-\binom b2\le a-2$.
It then follows that $a\ge2$. 
Note that $\beta=e(G[X])+e(G[Y])$. Then by Lemma~\ref{edge}, $e(G[Y])\ge \binom b2$, we have  \[
e(G[X])\le \binom a2+\binom b2+a-2-e(G[Y])\le \binom a2+a-2.
\]
Similarly, $e(G[Y])\le \binom b2+a-2\le \binom b2 +b-2$. 
By Lemma~\ref{lem:small-extremal-coloring},
$G[X]$ contains an $a$-clique $A$ and $G[Y]$ contains a $b$-clique $B$ such that deleting any vertex of the corresponding clique lowers the chromatic number  by at least one.

Since $|A|+|B|=a+b=\chi$ and $G$ contains no $K_\chi$, there are vertices $u\in A$ and $v\in B$ with
$uv\notin E(G)$. Color $G[X]-u$ with $a-1$ colors and $G[Y]-v$ with a disjoint set of $b-1$ colors,
and give $u$ and $v$ one new common color. This is a proper $(\chi-1)$-coloring of $G$, a contradiction.
Thus $\beta\ge B_\chi$. 
\end{proof}

Let $d_\chi=\binom{\chi}{2}-B_\chi$ and $L_\chi=\frac{d_\chi(d_\chi-1)}{B_\chi}$.

\begin{claim}\label{lem:critical-estimates}
There is an integer $a\ge2$ such that $s\ge\frac{(\chi-1)a+\chi-3}{2}$ and $\lambda_1\ge \chi-1+\frac{\chi-3}{\chi+a}$.
\end{claim}

\begin{proof}
Let $F$ be an induced $\chi$-vertex-critical subgraph of $G$. As $\omega(G)<\chi$, $F$ is not complete. By Lemma~\ref{order}, $|V(F)|\ge \chi+2$. Let $a=|V(F)|-\chi$. By  the fact that $e(F)\le m=\binom \chi2+s$ and Theorem~\ref{thm:dirac}, $2(\binom \chi2+s)\ge 2e(F)\ge(\chi-1)(\chi+a)+\chi-3$, so the first inequality follows. 
Moreover, by Lemma~\ref{subgraph} and Rayleigh's principle, we have
$\lambda_1(G)\ge\lambda_1(F)\ge\frac{2e(F)}{\chi+a}\ge \chi-1+\frac{\chi-3}{\chi+a}$,
which proves the second inequality.
\end{proof}

\noindent {\bf Case 2.2.1} $s\ge L_\chi$.


As $s\ge L_\chi$, we have $4B_\chi s-4d_\chi(d_\chi-1)\ge 0$, equivalently, 
\[
s+2d_\chi-1\le\sqrt{(s-1)^2+2\chi(\chi-1)s},
\] 
so $s+d_\chi\le U_\chi(s)$. It then follows by Lemma~\ref{maxcut} and Claim \ref{lem:maxcut-deficiency} that 
$\lambda^2\le m-\beta\le m-B_\chi=s+d_\chi\le U_\chi(s)$.

Suppose that $\lambda^2 = U_\chi(s)$. 
Then $s+d_\chi\le U_\chi(s)$, so  $s=L_\chi$.
If $\chi=2h$, then $B_\chi=h^2-1$, $d_\chi=h^2-h+1$ and $L_\chi=\frac{h(h^2-h+1)}{h+1}$. As $s$ is an integer,  $(h+1)|h(h^2-h+1)$. As $h(h^2-h+1)\equiv -3\pmod {h+1}$, we have $(h+1)|3$, which implies that $h=2$. So $\chi=4$ and $s=L_\chi=2$. However, we have by Claim~\ref{lem:critical-estimates}, $s\ge \frac{3a+1}{2}\ge \frac{7}{2}$, a contradiction. 
If $\chi=2h+1$, then 
$B_\chi=h^2+h-1$, $d_\chi=h^2+1$, and $L_\chi=\frac{h^2(h^2+1)}{h^2+h-1}$.
As $L_\chi$ is an integer, we have $(h^2+h-1)|h^2(h^2+1)$. Moreover, 
$h^2(h^2+1)\equiv3-4h\pmod{h^2+h-1}$. If $h\ge 3$, then
$(h^2+h-1)-(4h-3)=(h-1)(h-2)>0$, a contradiction. So $h=2$, $\chi=5$ and $s=L_\chi=4$.
By Claim~\ref{lem:critical-estimates} again, we have $s\ge \frac{4a+2}{2}\ge 5$, also a contradiction. So $\lambda^2< U_\chi(s)$.


\noindent {\bf Case 2.2.2.} $s< L_\chi$.



Suppose first that $\chi=4,5,6$. By a direct calculation, $L_4=2$, $L_5=4$ and $L_6=\frac{21}{4}$. However, Claim~\ref{lem:critical-estimates} gives $s\ge \frac{3\chi-5}{2}$, contradicting $s<L_\chi$. So $\chi\ge 7$.

Suppose that $\chi=7$. By Claim~\ref{lem:critical-estimates}, $s\ge8$. As $s$ is an integer and $s<L_7<9$, we have $s=8$. Again by Claim~\ref{lem:critical-estimates},
$8\ge\frac{6a+4}{2}=3a+2$, so $a=2$. Moreover, $\lambda_1(G)\ge6+\frac49=\frac{58}{9}$. So $\lambda^2\le2\left(\binom72+8\right)-\left(\frac{58}{9}\right)^2
 =\frac{1334}{81}<17$. On the other hand, $U_7(8)=\frac{9+\sqrt{721}}2>17$. Thus, $\lambda^2<U_7(8)$, as desired.


Suppose in the following that $\chi\ge8$.
If $\chi=2h$, then
$L_\chi=\frac{h(h^2-h+1)}{h+1}<h^2=\frac{\chi^2}{4}$.
If $\chi=2h+1$, then $L_\chi=\frac{h^2(h^2+1)}{h^2+h-1}$ and $\frac{\chi^2}{4}=h^2+h+\frac14$. As 
$$\left(h^2+h+\frac14\right)(h^2+h-1)-h^2(h^2+1)
 =\frac{8h^3-3h^2-3h-1}{4}>0,$$ we have $L_\chi<\frac{\chi^2}{4}$. So in both cases,  $L_\chi<\frac{\chi^2}{4}$. 
If $a\ge \chi$, then Claim~\ref{lem:critical-estimates} gives
$s\ge\frac{(\chi-1)\chi+\chi-3}{2}=\frac{\chi^2-3}{2}>\frac{\chi^2}{4}$,
contradicting $s<L_\chi<\frac{\chi^2}{4}$. Hence $a<\chi$, so $\chi+a\le2\chi-1$. It then follows by Claim~\ref{lem:critical-estimates} that
$\lambda_1(G)\ge \chi-1+\frac{\chi-3}{2\chi-1}\ge \chi-\frac23$. % as $\chi$\ge 8
So 
\[
\lambda^2\le 2m-\lambda_1^2\le \chi(\chi-1)+2s-\left(\chi-\frac23\right)^2
 =2s+\frac{\chi}{3}-\frac49. %=:V_k(s).
\]
We are left  to prove $2s+\frac{\chi}{3}-\frac49<U_\chi(s)$ whenever
$\frac{3\chi-5}{2}\le s<L_\chi$, which is equivalent to $3s+\frac{2\chi}{3}-\frac{17}{9}<\sqrt{(s-1)^2+2\chi(\chi-1)s}$. Since $s\ge \frac{3\chi-5}{2}$ and $\chi\ge 8$, we have $3s+\frac{2\chi}{3}-\frac{17}{9}\ge \frac{93\chi-169}{18}\ge 0$, so the above inequality is further equivalent to \[
324s^2-(81\chi^2-243\chi+378)s+18\chi^2-102\chi+104=:P_\chi(s)<0. 
\]
Note that
\begin{align*}
P_\chi\left(\frac{3\chi-5}{2}\right)&=-\frac{243\chi^3-2628\chi^2+7413\chi-6148}{2}\\
&=-\frac{243(\chi-8)^3+3204(\chi-8)^2+12021(\chi-8)+9380}{2}<0.
\end{align*} %$$P_\chi\left(\frac{3\chi-5}{2}\right)=-\frac{243\chi^3-2628\chi^2+7413\chi-6148}{2}<0.$$
Moreover, if $\chi=2h$, then $h\ge 4$ and
\begin{align*}
P_{\chi}(L_\chi) &=-\frac{2(3h^2-11h+13)(27h^3-12h^2+11h-4)}{(h+1)^2}\\
&=-\frac{2\left(3(h-4)^2+13(h-4)+17\right)\left(27(h-4)^3+312(h-4)^2+1211(h-4)+1576\right)}{(h+1)^2}\\
&<0,
\end{align*}
%$$P_{\chi}(L_\chi) =-\frac{2(3h^2-11h+13)(27h^3-12h^2+11h-4)}{(h+1)^2}<0,$$ 
and if $\chi=2h+1$, then $h\ge 4$, and  
\begin{align*}
P_\chi(L_\chi)&=-\frac{2(3h^3-14h^2+16h-10)(27h^4+15h^3+14h^2-7h+1)}{(h^2+h-1)^2}\\
&=-\frac{2\left(3(h-4)^3+22(h-4)^2+48(h-4)+22\right)}{(h^2+h-1)^2}\\
&\quad \cdot \left( 27(h-4)^4+447(h-4)^3+2786(h-4)^2+7737(h-4)+8069\right)\\
&<0.
\end{align*} 
As $P_\chi$ is convex in $s$, we have $P_\chi(s)<0$ whenever $\frac{3\chi-5}{2}\le s<L_\chi$. Therefore, $\lambda^2<U_\chi(s)$. %$\lambda^2\le V_{\chi}(s)<U_{\chi}(s)$.

Combining all cases, we complete the proof.
\end{proof}


\begin{thebibliography}{99}
\bibitem[BernshteynKostochka]{BernshteynKostochka} A. Bernshteyn, A. Kostochka, Sharp Dirac's theorem for DP-critical graphs, J. Graph Theory 88 (2018) 521--546.
\bibitem[CDS]{CDS} D.M. Cvetkovi\'{c}, M. Doob, H. Sachs, H. Spectra of Graphs, Johann Ambrosius Barth, Heidelberg, 1995.
\bibitem[Con]{Con} G. Constantine, Lower bounds on the spectra of symmetric matrices with nonnegative entries, Linear Algebra Appl. 65 (1985) 171--178.
\bibitem[Dirac]{Dirac} G.A. Dirac, A theorem of R. L. Brooks and a conjecture of H. Hadwiger, Proc. London Math. Soc. (3) 7 (1957) 161--195.
\bibitem[FYW]{FYW} Y.-Z. Fan, G.-D. Yu, Y. Wang, The chromatic number and the least eigenvalue of a graph, Electron. J. Combin. 19 (2012) \#P39.
\bibitem[HS]{HS} Y. Hong, J. Shu, Sharp lower bounds of the least eigenvalue of planar graphs, Linear Algebra Appl. 296 (1999) 227--232.
\bibitem[Jen]{Jen} T.R. Jensen, Dense critical and vertex-critical graphs, Discrete Math. 258 (2002) 63--84.
\bibitem[Pow]{Pow} D.L. Powers, Graph Partitioning by Eigenvectors, Linear Algebra Appl. 101 (1988) 121--133.
\bibitem[Powers]{Powers} D.L. Powers, Bounds on graph eigenvalues, Linear Algebra Appl. 117 (1989) 1--6.
\bibitem[TangElphick]{TangElphick} Q. Tang, C. Elphick, Proof of a conjectured spectral upper bound on the chromatic number of a graph, Electron. J. Combin. 33 (2026) \#P2.65.
\bibitem[WuElphick]{WuElphick} B. Wu, C. Elphick, Upper bounds for the achromatic and coloring numbers of a graph, Discrete Appl. Math. 217 (2017) 375--380.
\end{thebibliography}
\end{document}
